\section{Exact comparisons and succinct arithmetic}
\label{sec:exact-arithmetic}

Polynomial-time approximation at every requested accuracy does not
settle an exact order test. A nonzero quantity may require exponentially
many printed accuracy bits to distinguish from zero. For globally
strongly convex polynomial optimization, shared Newton circuits reach
that accuracy with polynomially many arithmetic operations. Algebraic
separation then converts exact comparisons, including equality, into
one $\PosSLP$ instance.

\subsection{One sign circuit for a polynomial observable}

\begin{theorem}[Strong-convexity upper reduction]
\label{thm:exact-upper}
Let $n\ge1$ and let $f,h\in\Q[X_1,\ldots,X_n]$ be explicitly supplied sparse
polynomials, with unary degree bounds. Suppose the input supplies
$\mu\in\Q$, $\mu>0$, and promises
\[
 \nabla^2f(X)\succeq\mu I\qquad(X\in\R^n).
\]
Let $p$ be the unique unconstrained minimizer of $f$. Each strict
order, weak order, and equality predicate on $h(p)$ versus zero has a
deterministic polynomial-time many-one reduction to one $\PosSLP$
instance. The circuit is constructed without $\PosSLP$ queries.
A rational threshold is absorbed in $h$ and counted in the input.
The same conclusion holds under any degree encoding guaranteeing
numerical degrees polynomially bounded by the input length.
\end{theorem}

The degree-four theorem and its observable extension are proved in the
\href{../../research-20260927/strong-convex-quartic-posslp-upper.md}{quartic upper-bound companion}.
The explicit sparse, unary-degree extension is proved in the
\href{../reviews/exact-core/general-degree/strong-convex-polynomial-posslp-upper.md}{variable-degree companion}.
The argument below states the common mechanism and its representation
requirements.

\paragraph{A polynomial-bit starting point.}
Let $L\ge2$ include the dimension, both polynomials, the degree
bounds, and $\mu$. Increasing it by a fixed constant if necessary,
assume $n,\deg f,\deg h\le L$, coefficient one-norms at most
$2^{2L}$, and $2^{-L}\le\mu\le2^L$. Strong convexity gives
coercivity and
\[
 \|p\|\le\|\nabla f(0)\|/\mu\le2^{3L}.
\]
Set $R=2^{4L}$ and $B=2^{16L^2}$. On $\|X\|\le R$,
$B$ bounds values and derivative operator norms through order three
for both polynomials. Indeed, an order-$j$ component is bounded by
$2^{2L}L^jR^L$, and passage to a tensor norm costs at most
$L^{j/2}$. The third derivative bound makes $B$ a Hessian Lipschitz
constant.

Put
\[
 \rho=\mu/(4B),\qquad \varepsilon_0=\mu^3/(32B^2).
\]
These rationals have polynomial binary length. Convex value
optimization gives a rational $x_0$ with
$f(x_0)-f(p)\le\varepsilon_0$ in polynomial bit time;
Corollary~1.2 of \citet{SlotSteurerWiedmer2025} applies directly. Strong convexity
implies $\|x_0-p\|\le\rho$. The radius-$\rho$ ball about $p$
is contained in the radius-$R$ ball used above. This preliminary
computation asks for polynomially many explicit accuracy bits and
does not use a sign oracle.

\paragraph{A separation bound used only in the proof.}
For $\alpha=h(p)$, the real solution set of
\begin{equation}
 \exists X:\quad\nabla f(X)=0,\qquad z=h(X)
 \label{eq:exact-singleton}
\end{equation}
is the singleton $\{\alpha\}$. Effective one-block real quantifier
elimination supplies a quantifier-free description whose polynomial
degrees and coefficient bit lengths are at most $2^{a(L)}$, for a
fixed effective polynomial $a$. This uses the degree estimate
$D^{O(n)}$ and $n\log D=\poly(L)$; the
primary theorem is \citet[Theorem~2.27]{Basu2014}.
The companion proof specifies its coefficient-height use.
The reduction does not execute quantifier elimination.

Some nonzero integer polynomial in that description must vanish at
$\alpha$; otherwise every sign condition would be locally constant
near $\alpha$, contradicting the singleton. Remove a power of $z$
from that polynomial. Its nonzero integer constant coefficient and
the elementary root bound give
\begin{equation}
 \alpha\ne0\quad\Longrightarrow\quad
 |\alpha|\ge g:=2^{-2^{a(L)+1}}.
 \label{eq:exact-separation}
\end{equation}
We enlarge $a$ so that $a(L)\ge32L^2+10$. The number $g$ has a
short circuit: start with $1/2$ and square $a(L)+1$ times. No
assumption on the dimension of the complex critical locus is needed;
\eqref{eq:exact-singleton} isolates one real value.

\paragraph{Newton iteration with shared arithmetic.}
Take exact rational Newton steps
\[
 x_{t+1}=x_t-\nabla^2f(x_t)^{-1}\nabla f(x_t).
\]
Within the established neighborhood, $e_t=\|x_t-p\|$ obeys
\[
 e_{t+1}\le\frac{B}{2\mu}e_t^2,\qquad
 e_t\le\frac{2\mu}{B}\,2^{-3\cdot2^t}.
\]
These estimates also keep every iterate in that neighborhood. A
positive definite Hessian is solved by rational
$LDL^{\mathsf T}$ elimination without pivoting, square roots, or
sign queries. Sparse differentiation and monomial evaluation add
polynomially many circuit gates per step. Shared nodes preserve the
values of previous steps without expanding their rational fractions.

After $k=a(L)+2$ steps, the bound on $\nabla h$ gives
\[
 |h(x_k)-\alpha|\le B e_k\le g/8.
\]
Let $\widehat\alpha=h(x_k)$. Each row of the following table gives
a rational circuit positive exactly for its listed predicate:
\begin{center}
\begin{tabular}{ll}
\toprule
Circuit output & Predicate\\
\midrule
$\widehat\alpha-g/2$ & $\alpha>0$\\
$\widehat\alpha+g/2$ & $\alpha\ge0$\\
$-\widehat\alpha-g/2$ & $\alpha<0$\\
$-\widehat\alpha+g/2$ & $\alpha\le0$\\
$g^2/4-\widehat\alpha^2$ & $\alpha=0$\\
\bottomrule
\end{tabular}
\end{center}
For zero, the approximation is within $g/8$; for a nonzero value,
its magnitude is at least $7g/8$. This verifies all equality cases
without an unknown nonzero-gap promise.

Represent every rational gate as $N/D$ with $D>0$. In particular,
division can be removed by
\[
 \frac{N_a/D_a}{N_b/D_b}
 =\frac{N_aD_bN_b}{D_aN_b^2}.
\]
All divisors are nonzero, by the positive-definite elimination
argument. Ordinary fraction rules handle the other operations, and
each rational gate adds only constantly many integer gates. Building
printed binary constants from $0$ and $1$ leaves a polynomial-size
integer circuit. Its final numerator is the required single
$\PosSLP$ instance. This proves Theorem~\ref{thm:exact-upper}.

The sparse input and degree bound are essential to the warm start and
the ordinary bit cost. This proof does not give a polynomial bound in
sparse binary-degree length alone or for a succinct circuit encoding
of the input polynomial. It also does not remove the supplied global
strong-convexity modulus.

\subsection{Checkable curvature and matching quartic lower bounds}

A useful input class replaces the semantic curvature promise with a
supplied rational identity
\begin{equation}
 v^{\mathsf T}\nabla^2f(X)v
 =w(X,v)^{\mathsf T}Mw(X,v),\qquad M\succ0,
 \label{eq:full-hessian-certificate}
\end{equation}
where the monomial vector $w$ is linear in $v$ and contains all
$v_i$. The matrix and monomial list are included in the input.
Sparse coefficient comparison checks the identity, and rational
symmetric elimination checks positive definiteness in polynomial
time. If $m=\dim M$, then
\[
 \mu=\det(M)/(\operatorname{tr}M)^{m-1}>0
\]
is a rational lower bound on $\lambda_{\min}(M)$ of polynomial
binary length. Because $\|w(X,v)\|^2\ge\|v\|^2$, it also
bounds the Hessian below. Enlarge $L$ to include this derived $\mu$;
the enlarged length is polynomial in the certificate input length.
For quartics we use the full basis
$w=(v,X\otimes v)$. Invalid certificates can be rejected before
applying the reduction, giving an ordinary decision language. The
construction does not recognize all strongly convex polynomials.

\begin{theorem}[Certified quartic order comparisons]
\label{thm:quartic-exact-complete}
For rational quartics on $\R^n$ supplied with a positive definite
rational Hessian Gram on $(v,X\otimes v)$, exact strict and weak
minimum-value comparisons and minimizer-coordinate comparisons are
$\PosSLP$-complete under polynomial-time many-one reductions.
The minimum-sign lower reduction has nonzero minimum on both source
branches. Equality has the upper bound of
Theorem~\ref{thm:exact-upper}; no matching equality lower bound is
asserted here.
\end{theorem}

The upper bound follows by choosing $h=f-r$ or $h=X_j-r$.
The lower reductions require a globally convex realization of
arithmetic computation, which is a separate substantive construction.
The
\href{../../research-20260927/posslp-certified-cubic-root-reduction.md}{signed-root sign compiler}
simulates a source integer circuit with bounded analytic signals,
and the
\href{../../research-20260927/signed-odd-root-circuit-quartic.md}{quartic realization}
places those signals at a unique zero of a rational sum of squares.
A carefully weighted exposing quadratic controls the otherwise
negative Hessian terms of squared residuals. The construction supplies
a positive definite Gram on the full Hessian basis, not merely a
positive semidefinite form evaluated on rank-one tensor vectors.

For minimum-sign comparison, a cubic perturbation encodes the signal
in the minimum without losing the full Hessian certificate. The
\href{../../research-20260927/unconstrained-quartic-posslp-reduction.md}{unconstrained sign reduction}
uses distinct signal coordinates $u,v$, with $u>0$, $v\ne0$,
$|v|\le1/8$, and $u^2\le|v|/2$ at the old zero. After scaling,
the perturbed objective has Hessian at least $I$. At that zero its
value is $-u^2v$ and its squared gradient norm is
$4u^2v^2+u^4$. If $v>0$, the value is already negative. If $v<0$,
strong convexity bounds the possible further decrease by half the
squared gradient norm, leaving a strictly positive minimum. The
scaling, the signal generator, and the full rational Gram all have
polynomial encoding length; they do not require printing the tiny
signals. The complete construction and its
\href{../reviews/exact-core/lower-review.md}{independent lower-bound review}
provide these quantitative interfaces.

The architecture of using accurate Newton circuits, algebraic
separation, and division elimination is established prior work
\citep{AllenderEtAl2009,EtessamiStewartYannakakis2012}. Exact semidefinite feasibility already has
arithmetic-circuit comparison reductions \citep{TarasovVyalyi2008}. The content
of Theorem~\ref{thm:quartic-exact-complete} is the classification for
the stated explicit quartic class and its strong certificates; it is
not a claim that exact convex feasibility or compressed Newton
simulation was first studied here.

\subsection{A rational optimizer and a known minimum do not settle coordinates}

\begin{theorem}[Rational-optimizer restriction]
\label{thm:rational-optimizer}
The coordinate-comparison lower bound in
Theorem~\ref{thm:quartic-exact-complete} holds even with all of the
following additional promises:
\begin{enumerate}[label=(\roman*)]
\item the unique optimizer $p$ lies in $\Q^N\cap[-1,1]^N$;
\item the minimum is the known value zero, and the input supplies
$N+1$ rational quadratic square factors;
\item $\nabla^2F\succeq(3/2)I$, with a supplied full positive
definite rational Hessian Gram;
\item the queried coordinate is nonzero and the threshold is zero.
\end{enumerate}
Rationality is a promise, not a property recognized by the theorem.
\end{theorem}

The
\href{../../research-20260927/rational-optimizer-posslp-coordinate-comparison.md}{complete reduction}
first compiles an integer source circuit with output $A$ into a shared
circuit of rational unit quaternions. Multiplication and inversion
preserve rationality and bounded coordinates. A designated coordinate
is nonzero with the sign of $2A-1$. A specialized quartic realization
has as its unique zero exactly the list of quaternion gate values,
with their coordinates preserved. It supplies the stated square
factors and full Hessian Gram in polynomial time. The
\href{../reviews/exact-core/rational-review.md}{independent review}
checks the sign compiler, realization, rational certificate, and
encoding bounds.

Every optimizer coordinate is rational, but expanding the gate
fractions is not part of the construction. Thus the result distinguishes
field membership from representation size and arithmetic comparison
cost. It concerns an optimizer coordinate. The minimum is already
known to be zero, and the theorem does not claim hardness of its sign
under this same promise. A perturbation encoding a minimum sign can
move the optimizer to an irrational point.

For the general upper theorem, the Newton vector $x_k$ is likewise a
rational vector represented by a polynomial-size shared circuit. When
$\min f<0$, choosing $h=f$ ensures $f(x_k)<0$, so the construction
supplies a short rational-circuit feasible witness without first
deciding feasibility. Checking its exact objective sign is still the
$\PosSLP$ comparison. This supplies neither polynomial-size expanded
fractions nor a rational feasible point when the minimum equals zero.

\subsection{Exact witnesses: existence and expanded size}

Even a strongly convex quartic with rational minimum need not have a
rational optimizer. The following small instance makes the distinction
between a point approximation and an exact feasible witness explicit.

\begin{proposition}[An irrational zero in two variables]
\label{prop:irrational-quartic-zero}
Put
\begin{align*}
 A(x,y)&=12599x^2-10000xy+7937y^2-15874x-12599y+20000,\\
 F(x,y)&=A(x,y)^2+10000\bigl((x^2-y)^2+(y^2-2x)^2\bigr).
\end{align*}
Then $F$ is a rational sum of squares with a supplied positive
definite rational full Hessian Gram, and
\[
 \nabla^2F\succeq4096I,\qquad
 \min_{\R^2}F=0,\qquad
 \argmin_{\R^2}F=\{(\sqrt[3]{2},\sqrt[3]{4})\}.
\]
Consequently $\{F\le0\}$ is nonempty and contains no rational
point. The same is true when the domain is $[1,2]^2$.
\end{proposition}

Any zero must satisfy $y=x^2$ and $y^2=2x$, hence
$x(x^3-2)=0$. The origin is excluded by $A(0,0)=20000$.
At $(r,r^2)$ with $r^3=2$, substitution gives $A=0$.
This proves the zero-set assertion. The curvature statement is an
exact rational identity, provided in the
\href{../../research-20260927/convex-quartic-rational-sos.md}{full Hessian certificate}:
after an invertible rational affine change, the Hessian biform minus
$4096\|v\|^2$ is represented by a strictly positive definite
rational matrix on all six monomials
$(v_1,v_2,xv_1,xv_2,yv_1,yv_2)$. Its displayed diagonal-dominance
certificate and invertible congruence prove positivity; numerical
matrix tests are not part of the argument. The
\href{../reviews/exact-core/singleton-review.md}{independent singleton audit}
checks this identity and the zero-set proof.

The example answers negatively the rational-feasible-point witness
question for quartics in Table~1 of \citet{SlotSteurerWiedmer2025}, version~1.
It is consistent with that paper's univariate result: the present
quartic is bivariate. Its optimizer has a short algebraic description,
so the absence of a rational point excludes neither other certificate
formats nor NP membership of the decision problem. A full positive
definite \emph{ordinary polynomial} Gram cannot represent this $F$,
because it vanishes at a real point and the ordinary monomial vector
contains the constant one. Its positive definite \emph{Hessian} Gram
represents a different polynomial in a different basis.

\begin{theorem}[Expanded witness bounds]
\label{thm:expanded-witnesses}
The following assertions concern ordinary binary fractions.
\begin{enumerate}[label=(\roman*)]
\item There are polynomial-size rational quartics in $n=2k$
variables, with supplied full Hessian Grams at least the identity,
whose zero sublevel sets are bounded and have nonempty interior,
but every rational feasible point requires
$\Omega(n2^{n/2})$ bits. All feasible coordinate magnitudes are
bounded by an absolute constant.
\item There are strictly positive rational quartics with such
Hessian certificates that admit both a short rational positive
semidefinite ordinary Gram and a rational positive definite ordinary
Gram. Every rational positive definite ordinary Gram requires
$\Omega(n2^{n/2})$ total denominator bits. The basis is the
full vector of monomials of degree at most two, each occurring once.
\item Every strictly positive rational quartic supplied with a
positive definite full rational Hessian Gram has a positive definite
rational ordinary Gram with total expanded size
$\poly(L)2^{O(n)}$, where $L$ counts the polynomial and supplied
Hessian Gram. This is a size bound, not a polynomial-time construction
in $L$ alone.
\end{enumerate}
\end{theorem}

The complete proofs are the
\href{../../research-20260927/strict-convex-quartic-rational-witness-lower-bound.md}{feasible-point lower bound},
\href{../../research-20260927/interior-gram-bit-lower-bound.md}{positive definite Gram lower bound},
and
\href{../../research-20260927/interior-gram-single-exponential-upper.md}{Gram size upper bound}.
The lower constructions encode a positive signal at most
$1000^{-(k+3)2^k}$ using $k$ bounded cubic-root gates. One
perturbation confines the whole feasible body extremely close to a
cubic irrational coordinate; the integer difference of cubes then
forces a large rational denominator. Another perturbation produces a
strictly positive polynomial with an extremely small value. A uniform
trace bound for its ordinary Grams and that small Rayleigh quotient
force a tiny positive determinant, hence long denominator encodings.
This second argument requires positive definiteness; the same family
has a short singular Gram.

For the upper bound, rational sampling in a nonempty open
semialgebraic set supplies a center of singly exponential size where
$f-\|\nabla f\|^2/(2\mu)>0$. Taylor integration of the supplied
Hessian Gram then gives an ordinary positive definite rational Gram
with the claimed size. The
\href{../reviews/exact-core/witness-review.md}{witness audit}
checks the sampling theorem, the full-span Taylor construction, and
both lower bounds. Exponential dependence on dimension is therefore
qualitatively necessary and sufficient for these interior Grams;
the exponent constants are not matched. The lower bounds are
superpolynomial in the constructed input length, but are not asserted
as $2^{\Omega(L)}$. None excludes a compact arithmetic or algebraic
description. Section~\ref{sec:certificates} states the corresponding
certificate interfaces.

Finally, a globally convex polynomial of degree at most three is
quadratic: its affine Hessian must remain positive semidefinite in
both directions of every line, forcing each nonconstant matrix
coefficient to vanish. Strongly convex rational quadratic minimization
on all of $\R^n$ reduces to rational linear algebra. Degree four is
therefore the first degree supporting the constructions above. This
degree statement and $\PosSLP$-completeness do not prove a separation
from P, an NP-hardness theorem, or a practical exact solver. The
\href{../literature/exact-prior.md}{primary-source comparison}
records the established methodological precedents and the remaining
publication-priority limits.
